\section{The Institutional Response: The Anthropic/METR Arrangement}

On 12 September 2026, Anthropic CEO Dario Amodei published ``Pace the Frontier,'' a three-part essay calling for embedded third-party evaluators and a coordinated industry slowdown~\cite{amodei2026,bbc2026,techcrunch2026}. Sam Altman and Elon Musk publicly endorsed the call within hours~\cite{techcrunch2026}. The timing was coordinated and the substance was convergent. The statements do not, on their own, establish motive or effect. What they establish is that the leading laboratories converged, within a short window, on a specific institutional form: the embedded evaluator with privileged access. (The essay characterizes the endorsement as coming from the CEOs of Anthropic, OpenAI, xAI, and Google DeepMind; only the first three could be independently verified against primary reporting. The fourth is reported in the secondary literature but is not confirmed by the primary sources reviewed here. This caveat applies throughout Section~4.)

\subsection{What the documents show}

The most documented instantiation is the Anthropic/METR arrangement, announced in October 2026~\cite{anthropic2026a}. Six METR employees are embedded at Anthropic with company hardware, office space, and access ``roughly on par'' with internal risk teams. METR can publish its risk assessments; Anthropic cannot veto publication but can require redaction of predefined categories: customer data, attorney-client privileged information, and information about model architectures. Three access tiers are being built: Tier~1 (narrow scope, Anthropic approves publication), Tier~2 (broader access, publication without approval but Anthropic retains redaction rights), and Tier~3 (employee-like access, the METR arrangement). Accenture is a parallel evaluator under a partnership announced 18 September 2026 in which both firms committed approximately \$1 billion each over five years to embedded-evaluation work; Anthropic funds Accenture directly~\cite{anthropic2026b,cnbc2026,wsj2026}. METR is not paid by Anthropic; it is philanthropically and government funded. (An earlier draft of this essay referred to Faculty as a third parallel evaluator; the reference could not be confirmed against a primary source and has been removed pending verification.)

The arrangement's design features matter more than its stated purpose. Measured against the eight questions the editorial framework requires:

\begin{itemize}
\item \textbf{Who is subject:} Anthropic only, voluntarily. The draft bill would extend to ``frontier developers'' defined by capability thresholds.
\item \textbf{What triggers obligations:} Anthropic's own decision to invite METR. The draft bill triggers on release of a frontier model above a capability threshold.
\item \textbf{What costs are imposed:} METR estimates that a rigorous two-month investigation requires API credits in the millions of dollars~\cite{metr2026a}; the precise figure has not been confirmed against a primary source and is treated here as order-of-magnitude rather than exact. METR is not paid by Anthropic; Accenture is. The asymmetry is the point: the deepest-access evaluator is the one not paid by the firm.
\item \textbf{Who sets standards and controls evaluation:} Anthropic defines access tiers, redaction categories, and Tier~1 scope. METR has publication rights but no veto over release. The draft bill would have Commerce register evaluators but not set standards.
\item \textbf{Proportionality:} Unknown. The tier structure implies proportionality, but the criteria for assigning tiers are not public.
\item \textbf{Shared infrastructure, exemptions, simplified pathways:} None. The model is bespoke: each lab negotiates its own arrangement. The draft bill's Commerce registry is a coordination mechanism, not a shared facility.
\item \textbf{Differential effects:} Established labs can absorb API costs in the millions, have legal and security infrastructure, and can host embedded teams. Smaller developers cannot. Open-weight developers are excluded by design---there is no ``employee-like access'' to a model whose weights are public.
\item \textbf{Less-restrictive alternatives:} Government-funded evaluation (Anthropic's stated preference, not yet realized); shared public evaluation infrastructure (no proposal); standardized disclosure without embedded access (weaker but cheaper); open-weight-specific regimes such as post-hoc audit.
\end{itemize}

The arrangement is therefore \textbf{contestable on publication}---METR has published critical findings, including the conclusion that Anthropic's monitoring systems contained ``many vulnerabilities'' exploitable by misaligned agents~\cite{metr2026b}---and \textbf{potentially exclusionary on access, proportionality, and open-weight treatment}. The word ``potentially'' is load-bearing: lack of shared infrastructure is an observed feature; its exclusionary effect is a further claim that requires evidence about actual or expected entry costs, feasible alternative arrangements, and the conditions under which smaller developers would be subject to the proposed regime. That evidence is not yet available. The design audit identifies where the burden of proof lies; it does not discharge it.

The draft bill Anthropic has circulated would generalize the template. According to reporting by Politico (via~\cite{wifc2026}) and Inside AI Policy~\cite{insideaipolicy2026}, the bill requires frontier developers to submit to independent evaluation, requires evaluators to register with the Department of Commerce, allows labs to pay evaluators provided compensation is not tied to results, and overrides certain state laws such as California's. The bill is reported to permit government intervention only when a model poses a foreseeable, imminent risk of more than \$1 billion in harm~\cite{wifc2026}. A federal standard designed by the firm displaces a patchwork of state standards that might have been more demanding. The compliance cost structure favors incumbents: evaluator registration, API costs, legal and administrative overhead. Whether the codified regime is contestable or exclusionary depends on design choices not yet fixed. This essay's analysis of the draft bill rests on secondary reporting (Politico, Inside AI Policy) because the bill text itself has not been publicly released; the analysis should be revisited when the bill is introduced and the text is available.

\subsection{The strategic-incentive hypothesis}

The institutional response described above is consistent with a strategic-incentive hypothesis: as the economic gradient flattens, frontier firms may have an increased incentive to favor governance designs that perform the maintenance function while preserving their position within the infrastructure layer. This is a hypothesis, not a finding, and the paper does not establish that the incentive is stronger now than in other periods. The incentive to shape regulation may be persistent across the industry; the deployment transition may amplify it, or may simply coincide with it.

Three non-exclusive explanations deserve equal consideration.

\textbf{H1 --- Strategic incentive.} The proposals are a response to the flattening gradient: firms seek to shape the configuration that performs the maintenance function, favoring parameters that preserve their position in the infrastructure layer.

\textbf{H2 --- Genuine safety concern.} The proposals are a response to documented incidents---the Hugging Face breach, the false homicide tip, the FTC's civil investigative demands---and reflect a good-faith assessment that external oversight is necessary. The timing reflects the incidents, not the gradient.

\textbf{H3 --- Regulatory uncertainty and legal certainty.} The proposals are a response to the patchwork of state regulations and the threat of liability. Firms prefer a federal standard they can shape to a state-by-state regime they cannot control, regardless of whether the standard is exclusionary.

These explanations are compatible, not mutually exclusive. All three may be operating simultaneously. The evidence reviewed here cannot adjudicate among them, and the paper does not attempt to. What the paper can establish is narrower: that the design features of the response have distributional consequences independently of which explanation is correct, and that the framework for evaluating those features is available now, before codification fixes them.

H1 contains an assumption that should be exposed. It assumes that firms have an incentive to favor configurations that preserve their position within the system. That is plausible, but the paper has not shown that flattening the economic gradient increases this incentive relative to other periods. The incentive to shape regulation may be persistent across the industry, not something newly caused by the deployment transition. The hypothesis is therefore stated with a directional qualifier---may have an increased incentive---rather than as a causal claim. What distinguishes H1 from H2 and H3 is not the presence of the incentive but its predicted design signature: H1 predicts that the proposals will favor parameters that raise compliance costs for entrants relative to incumbents, even when the safety benefit is marginal. H2 predicts that the proposals will track the severity of documented incidents. H3 predicts that the proposals will favor federal preemption of state law, regardless of their effect on entry. The Anthropic draft bill is consistent with all three: it requires independent evaluation (H2), overrides California law (H3), and imposes fixed costs that fall disproportionately on entrants (H1). The paper's contribution is not to choose among them, but to show that the design features are consequential regardless of which explanation is correct.

The design audit in Section~5 does not depend on H1 being true. It evaluates the arrangement's design features and their distributional effects, whatever the motive. A regime designed with the best of intentions can produce exclusionary effects; a regime designed with cynical motives can produce genuine safety benefits. What matters is what the design does---who is subject to it, who controls it, who bears its costs, who can contest its findings.
